% ECONOMICS_PROBLEM: {"id": "H21-58", "title": "Coordinating multiple tax instruments", "jel_code": "H21", "source_id": "OP-0471"}
% !TeX program = xelatex
\documentclass[11pt,a4paper]{article}
\usepackage[margin=23mm]{geometry}
\usepackage{fontspec}
\setmainfont{Latin Modern Roman}
\usepackage{amsmath,amssymb,mathtools}
\usepackage{xcolor,enumitem,booktabs,longtable,array,xurl,hyperref}
\definecolor{muted}{HTML}{53636C}
\hypersetup{colorlinks=true,urlcolor=blue,linkcolor=blue}
\setlength{\parindent}{0pt}
\setlength{\parskip}{5pt}
\newcommand{\R}{\mathbb R}
\newcommand{\E}{\mathbb E}
\newcommand{\N}{\mathbb N}
\newcommand{\ind}{\mathbf 1}
\newcommand{\Var}{\operatorname{Var}}
\newcommand{\Cov}{\operatorname{Cov}}
\newcommand{\supp}{\operatorname{supp}}
\DeclareMathOperator*{\argmin}{arg\,min}
\DeclareMathOperator*{\argmax}{arg\,max}
\newcommand{\entrymeta}[3]{{\sffamily\small\color{muted}\textbf{\texttt{#1}}\quad|\quad #2\par #3\par}}
\newcommand{\Problemref}[1]{\texttt{#1}}
\newcommand{\JELref}[2]{\texttt{#2} (JEL \texttt{#1})}
\begin{document}
\section*{Coordinating multiple tax instruments}
\textbf{Problem ID:} \texttt{H21-58}\quad \textbf{JEL:} \texttt{H21}\par
% BEGIN ECONOMICS STATEMENT
\label{problem:OP-0471}
{\sffamily\small\color{muted}Original subject group: Taxation and social insurance\par}
\label{definitions:problem:30007753}
\textbf{Audit classification:} Published research agenda. The integrating-instruments agenda is verified. Added a finite-horizon equilibrium definition, common numeraire, instrument equivalence criterion, and supremum/attainment distinction.
\subsection*{Definitions and precise research target}
Fix a finite horizon \(H\), a population of private household types \(\theta\sim F\), observable characteristics \(x(\theta)\), production and trading technologies, initial assets, an information structure and terminal conditions. A policy \(p=(T_y,T_k,t_c,T_f)\in\mathcal P\) specifies the complete labor- and capital-income schedules, commodity taxes and family transfers at each date. The admissible set \(\mathcal P\), administrative resource cost \(C(p)\), tax bases, remittance dates and any behavioral misperceptions are declared in advance. Every schedule uses only legally observable variables. A competitive equilibrium \(e\in\mathcal E(p)\) comprises household plans, firm plans, prices, receipts and market-clearing conditions at every date, with sequential budget feasibility and the specified terminal conditions. Exclude policies without an equilibrium and assume bounded integrable utilities.

Let \(U_\theta(e)\) be the declared lifetime welfare index, \(G_\theta\) the fixed social valuation, and \(B(p,e)\) the discounted receipts net of transfers and administration. Discount all public flows using the specified date-zero numeraire. For funding requirement \(R\), define
\[
 V_{\mathcal P}=\sup_{p\in\mathcal P}\inf_{e\in\mathcal E(p)}\int G_\theta(U_\theta(e))\,dF(\theta),
 \qquad B(p,e)\geq R\quad\forall e\in\mathcal E(p).
\]
The supremum ranges only over policies satisfying these conditions. A baseline comparison uses either a declared equilibrium selection or the same worst-equilibrium criterion throughout. Characterize joint reforms whose welfare gain is positive and whose revenue constraint holds under that criterion; characterize when a smaller instrument class \(\mathcal P_0\subseteq\mathcal P\) has value \(V_{\mathcal P_0}=V_{\mathcal P}\). The associated welfare loss is \(V_{\mathcal P}-V_{\mathcal P_0}\), not a difference between arbitrarily selected equilibria. Policies are allocation-equivalent only if their equilibrium allocations and net public receipts coincide under the specified model; matching one labor wedge does not establish equivalence. Existence of an optimizer requires additional compactness and continuity arguments. This is a newly scoped policy problem.
\subsection*{Variants and assumption changes}
\begin{enumerate}
\item \textbf{Static separability (editorial benchmark).} Use one period, identical homothetic consumption preferences weakly separable from labor, perfect observability and unrestricted nonlinear earnings taxes. Test instrument redundancy under these explicit restrictions, rather than assuming redundancy in a heterogeneous dynamic economy.
\item \textbf{Restricted instruments (editorial).} Limit earnings taxes to finitely many brackets and add commodity taxes. Compute \(V_{\mathcal P}-V_{\mathcal P_0}\) when the smaller class has uniform commodity taxation; this differs from unrestricted income taxation.
\item \textbf{Behavioral welfare (editorial).} Specify distinct decision and welfare utilities and an observation model for misperceived taxes. The optimizer then corrects distortions as well as redistributing resources.
\end{enumerate}
\subsection*{References and source audit}
Kaplow's 2022 manuscript section 7 provides the integration framework; section 9 discusses extensions. Journal metadata are verified on the 2024 AEA page. The exact multi-instrument minimax question is editorial.
\begin{enumerate}\raggedright
\item\hypertarget{definitions:source:30007753:1}{} Louis Kaplow. 2022. \emph{Optimal Income Taxation}. Olin Discussion Paper 1083, July 2022; sections 3.2, 7, 8.2 and 9, printed pp. 89–90.
\url{https://laweconcenter.law.harvard.edu/wp-content/uploads/2024/11/Kaplow_1083.pdf}.
DOI: \url{https://doi.org/10.3386/w30199}.
\item\hypertarget{definitions:source:30007753:2}{} Louis Kaplow. 2024. \emph{Optimal Income Taxation}. Journal of Economic Literature 62(2), 637–738; publisher abstract.
\url{https://www.aeaweb.org/articles?id=10.1257/jel.20221647}.
DOI: \url{https://doi.org/10.1257/jel.20221647}.
\end{enumerate}

\subsection*{Common conventions: Source mathematical conventions}

These conventions apply to each entry unless explicitly replaced.

\textbf{Models and identification.} A primitive model \(m\) specifies all probability laws, feasible choices, information, equilibrium and measurement rules used by its target. The admissible class \(\mathcal M\) is fixed independently of outcomes. If \(P_O(m)\) is its observable probability law and \(\tau(m)\) the target, its identified set at observable law \(P\) is
\[
 \mathcal I_\tau(P)=\{\tau(m):m\in\mathcal M,\ P_O(m)=P\}.
\]
Point identification means that this set is a singleton. An empty set signals incompatibility of data and assumptions. Coordinate bounds need not describe the joint set. Bounds are sharp only if every retained target is attainable or arbitrarily approachable by compatible models. Finite bounds require explicit restrictions on unobserved quantities; unrestricted confounding, hidden wealth, extrapolation or arbitrary equilibrium selection can yield uninformative sets.

\textbf{Causal interventions.} A potential outcome \(Y_i(p)\) is the outcome under a complete policy regime \(p\), including timing, financing, information and market spillovers. The target population and units are fixed before treatment. With interference, use the complete assignment vector or a justified exposure mapping; individual no-interference notation alone cannot identify an economy-wide intervention. A difference of mean potential outcomes does not require the cross-world joint distribution. Conditioning on post-treatment migration, survival, enrollment or employment changes the estimand and needs separate justification.

\textbf{Statistical statements.} An estimator, confidence set, forecast or test specifies a sampling law, model class and loss/coverage criterion. Uniform \((1-\alpha)\)-coverage means \(\inf_{m\in\mathcal M}\Pr_m\{\tau(m)\in C_n\}\geq1-\alpha\), or an explicitly asymptotic version. The whole parameter space trivially has coverage, so informativeness requires a stated width, risk or power objective. A forecasting improvement is relative to a named benchmark, information set and evaluation population.

\textbf{Equilibrium and optimization.} An equilibrium comprises feasible plans, optimality, beliefs where needed, budget constraints and all market/resource clearing. The primitives or a stated benchmark define each correspondence \(\mathcal E(p;m)\). Nonemptiness is assumed or proved, never inferred just from compact policy sets. A selected-equilibrium target uses a declared selection rule; a robust target quantifies over all permitted equilibria. Use suprema or infima unless attainment is established. An \(\varepsilon\)-optimal policy is feasible and within \(\varepsilon\) of the finite optimum value. Report an empty feasible set separately.

\textbf{Welfare and accounting.} Normative utility, interpersonal normalization, social weights, numeraire, discounting and the use of public receipts are primitives. Behavioral utility need not coincide with the welfare criterion. Household welfare includes ownership income and rents once; adding profits again to compensating variation generally double-counts them. A monetary equivalent is defined by an explicit transfer and price convention, with monotonicity and range conditions. Average effects, mechanism contributions, transfers and welfare losses are different targets. Infinite sums require convergence and dynamic feasible plans require terminal or no-Ponzi conditions.

\textbf{Source versions.} A working-paper date, revision date and journal publication year are distinguished. A DOI for a journal version is not automatically the DOI of an earlier manuscript. Locators refer to the stated version and printed pages where specified. A metadata-only check does not verify a claim about a full-text conclusion. Every entry's source subsection narrows the provenance claim accordingly.
% END ECONOMICS STATEMENT
\end{document}
